% !TEX root = ../main.tex
%% Chapitre 1 — Suites numériques
%% Sources : frontend/cours/resume/terminal/resume_suites_numeriques_terminal.txt
%%           frontend/cours/resume/terminal/resume_limites_convergence_suites_terminal.txt
%%           frontend/exercices/terminal_spé_Maths/exercice_limites_convergence_suites.txt
%% Exercices 13 à 15 : problèmes inédits rédigés pour le livre.

\begin{ouverture}[suites]{1}{Suites numériques}{Partie 1 · Analyse}
  {Récurrence\enspace·\enspace Limites\enspace·\enspace Comparaison\enspace·\enspace
   Convergence monotone\enspace·\enspace Suites récurrentes}
  \objectif{Démontrer une propriété par récurrence.}
  \objectif{Déterminer la limite d'une suite : opérations, factorisation, théorèmes de comparaison.}
  \objectif{Utiliser le théorème des gendarmes et le théorème de convergence monotone.}
  \objectif{Étudier une suite définie par $u_{n+1}=f(u_n)$ : encadrement, monotonie, limite.}
  \objectif{Exprimer une suite arithmético-géométrique en fonction de $n$ et programmer un algorithme de seuil.}
  \prerequis{Suites arithmétiques et géométriques (Première).}
  \prerequis{Sens de variation d'une fonction à l'aide de sa dérivée.}
  \prerequis{Second degré : factorisation et signe d'un trinôme.}
\end{ouverture}

%% =====================================================================
%%  L'ESSENTIEL
%% =====================================================================
\essentiel

\essSection{1}{Suites arithmétiques et géométriques}

{\renewcommand{\arraystretch}{1.55}\arrayrulecolor{otRule}%
\begin{tabularx}{\linewidth}{@{}>{\otSemi\fontsize{7.5}{10}\selectfont\leavevmode\color{otNavy}}p{27mm}YY@{}}
\toprule
 & {\otSemi\fontsize{9}{11}\selectfont Arithmétique de raison $r$}
 & {\otSemi\fontsize{9}{11}\selectfont Géométrique de raison $q$} \\
\midrule
\otCaps[8]{Définition} & $u_{n+1}=u_n+r$ & $u_{n+1}=q\,u_n$ \\
\otCaps[8]{Terme général} & $u_n=u_0+nr$\quad et\quad $u_n=u_p+(n-p)r$
  & $u_n=u_0\,q^{n}$\quad et\quad $u_n=u_p\,q^{n-p}$ \\
\otCaps[8]{Somme} & $\displaystyle u_0+\dots+u_n=(n+1)\,\frac{u_0+u_n}{2}$
  & $\displaystyle u_0+\dots+u_n=u_0\,\frac{1-q^{n+1}}{1-q}$\enspace($q\neq1$) \\
\otCaps[8]{Variations} & croissante si $r>0$, décroissante si $r<0$
  & si $u_0>0$ : croissante si $q>1$, décroissante si $0<q<1$ \\
\bottomrule
\end{tabularx}}

\vspace{2mm}
Deux sommes à connaître :\quad
$\displaystyle 1+2+\dots+n=\frac{n(n+1)}{2}$\qquad et\qquad
$\displaystyle 1+q+\dots+q^{n}=\frac{1-q^{n+1}}{1-q}$\enspace pour $q\neq1$.

\begin{savoirfaire}[Suite arithmético-géométrique $u_{n+1}=a\,u_n+b$, avec $a\neq1$]
\begin{enumerate}[label=\otLabel{\arabic*.},leftmargin=1.6em,itemsep=0.6mm,topsep=0.5mm]
  \item On cherche le réel $\ell$ tel que $\ell=a\ell+b$, soit $\ell=\dfrac{b}{1-a}$.
  \item On pose $v_n=u_n-\ell$ : alors $v_{n+1}=a\,v_n$, donc $(v_n)$ est géométrique de raison $a$.
  \item On en déduit $v_n=(u_0-\ell)\,a^{n}$, puis $u_n=\ell+(u_0-\ell)\,a^{n}$.
\end{enumerate}
\end{savoirfaire}

\essSection{2}{Le raisonnement par récurrence}

\begin{theoreme}[Principe de récurrence]
Soit $P(n)$ une propriété qui dépend d'un entier naturel $n$, et $n_0$ un entier naturel. Si :
\begin{itemize}[label={\color{otNavy}\textbullet},leftmargin=1.4em,itemsep=0.6mm,topsep=0.8mm]
  \item \etape{Initialisation} $P(n_0)$ est vraie ;
  \item \etape{Hérédité} pour tout entier $k\geq n_0$, si $P(k)$ est vraie, alors $P(k+1)$ est vraie ;
\end{itemize}
alors $P(n)$ est vraie pour tout entier $n\geq n_0$.
\end{theoreme}

\begin{attention}
Dans l'hérédité, on suppose $P(k)$ vraie pour \emph{un} entier $k$ fixé (c'est l'hypothèse de
récurrence), et on doit s'en servir pour démontrer $P(k+1)$. L'initialisation est indispensable :
la propriété « $4^{n}+1$ est divisible par $3$ » est héréditaire, puisque
$4^{k+1}+1=4\,(4^{k}+1)-3$, et pourtant elle est fausse pour tout entier $n$.
\end{attention}

\essSection{3}{Limite d'une suite}

\begin{definition}[Suite convergente]
La suite $(u_n)$ \textbf{converge} vers le réel $\ell$ si tout intervalle ouvert contenant $\ell$
contient tous les termes $u_n$ à partir d'un certain rang. On note $\displaystyle\lim_{n\to+\infty}u_n=\ell$.
Cette limite, si elle existe, est unique. Une suite qui ne converge pas est dite \textbf{divergente}.
\end{definition}

\begin{center}
\begin{tikzpicture}
\begin{axis}[otSchema, width=128mm, height=46mm,
  xmin=0, xmax=16.5, ymin=0, ymax=3.6, xtick=\empty, ytick=\empty,
  xlabel={$n$}, xlabel style={at={(1,0)},anchor=west}]
  \fill[otGreen,opacity=0.13] (axis cs:0,1.65) rectangle (axis cs:16.5,2.35);
  \draw[otGreen,line width=0.5pt] (axis cs:0,1.65) -- (axis cs:16.5,1.65);
  \draw[otGreen,line width=0.5pt] (axis cs:0,2.35) -- (axis cs:16.5,2.35);
  \draw[otNavy,dashed,line width=0.6pt] (axis cs:0,2) -- (axis cs:16.5,2);
  \draw[otMuted,dashed,line width=0.5pt] (axis cs:4.6,0) -- (axis cs:4.6,3.4);
  \addplot[only marks,mark=*,mark size=1.6pt,draw=otNavy,fill=white,
    samples at={0,1,2,3,4}] {2+1.4*cos(180*x)*0.75^x};
  \addplot[only marks,mark=*,mark size=1.6pt,otGreen,
    samples at={5,6,...,16}] {2+1.4*cos(180*x)*0.75^x};
  \node[anchor=east,font=\small,text=otNavy] at (axis cs:0,2) {$\ell$};
  \node[anchor=east,font=\scriptsize,text=otGreenDark] at (axis cs:0,2.42) {$\ell+\varepsilon$};
  \node[anchor=east,font=\scriptsize,text=otGreenDark] at (axis cs:0,1.58) {$\ell-\varepsilon$};
  \node[anchor=north,font=\otMedium\fontsize{7}{8}\selectfont,text=otMuted] at (axis cs:4.6,0) {rang $N$};
  \node[anchor=south west,align=left,font=\otLight\fontsize{7.5}{9.5}\selectfont,text=otInk]
    at (axis cs:7.2,2.6) {À partir du rang $N$, tous les termes\\sont dans l'intervalle $\intOO{\ell-\varepsilon}{\ell+\varepsilon}$.};
\end{axis}
\end{tikzpicture}
\end{center}

\begin{definition}[Limite infinie]
$\displaystyle\lim_{n\to+\infty}u_n=+\infty$ si tout intervalle $\intOO{A}{+\infty}$ contient tous les termes $u_n$
à partir d'un certain rang. On définit de même $\displaystyle\lim_{n\to+\infty}u_n=-\infty$.
\end{definition}

\noindent
\begin{minipage}[t]{0.47\linewidth}
\otIntertitre{Limites de référence}
{\renewcommand{\arraystretch}{1.5}\arrayrulecolor{otRule}%
\begin{tabular}{@{}l@{\hspace{8mm}}l@{}}
\toprule
$\dfrac{1}{n}$,\enspace $\dfrac{1}{n^{2}}$,\enspace $\dfrac{1}{\sqrt{n}}$ & tendent vers $0$ \\
$n$,\enspace $n^{2}$,\enspace $\sqrt{n}$ & tendent vers $+\infty$ \\
\bottomrule
\end{tabular}}
\end{minipage}\hfill
\begin{minipage}[t]{0.47\linewidth}
\otIntertitre{Limite de $(q^n)$}
{\renewcommand{\arraystretch}{1.5}\arrayrulecolor{otRule}%
\begin{tabular}{@{}l@{\hspace{8mm}}l@{}}
\toprule
$q>1$ & $\displaystyle\lim_{n\to+\infty}q^{n}=+\infty$ \\
$q=1$ & $q^{n}=1$ pour tout $n$ \\
$-1<q<1$ & $\displaystyle\lim_{n\to+\infty}q^{n}=0$ \\
$q\leq-1$ & pas de limite \\
\bottomrule
\end{tabular}}
\end{minipage}

\vspace{4mm}
\otIntertitre{Opérations sur les limites}
Les limites de sommes, produits et quotients se calculent « naturellement », sauf dans quatre cas
appelés \textbf{formes indéterminées} :
\[
  \infty-\infty \qquad 0\times\infty \qquad \frac{\infty}{\infty} \qquad \frac{0}{0}
\]

\begin{savoirfaire}[Lever une forme indéterminée]
\begin{itemize}[label={\color{otGreenDark}\textbullet},leftmargin=1.4em,itemsep=0.6mm,topsep=0.5mm]
  \item Polynôme ou quotient de polynômes : factoriser par le terme de plus haut degré.
  \item Différence de racines carrées : multiplier et diviser par la quantité conjuguée.
  \item Expression simplifiable : simplifier \emph{avant} de passer à la limite.
\end{itemize}
\end{savoirfaire}

\essSection{4}{Comparaison et convergence}

\begin{theoreme}[Comparaison]
Si, à partir d'un certain rang, $u_n\leq v_n$ et $\displaystyle\lim_{n\to+\infty}u_n=+\infty$, alors
$\displaystyle\lim_{n\to+\infty}v_n=+\infty$. De même, si $v_n\leq u_n$ et
$\displaystyle\lim_{n\to+\infty}u_n=-\infty$, alors $\displaystyle\lim_{n\to+\infty}v_n=-\infty$.
\end{theoreme}

\begin{theoreme}[Théorème des gendarmes]
Si, à partir d'un certain rang, $u_n\leq v_n\leq w_n$, et si $(u_n)$ et $(w_n)$ convergent vers
\emph{la même} limite $\ell$, alors $(v_n)$ converge vers $\ell$.
\end{theoreme}

\noindent
\begin{minipage}[c]{0.5\linewidth}
\otIntertitre{Exemple}
Pour tout $n\geq1$, $-1\leq\sin n\leq1$, donc
\[ -\frac{1}{n}\leq\frac{\sin n}{n}\leq\frac{1}{n}. \]
Les deux suites qui encadrent tendent vers $0$, donc, d'après le théorème des gendarmes,
$\displaystyle\lim_{n\to+\infty}\frac{\sin n}{n}=0$.
\end{minipage}\hfill
\begin{minipage}[c]{0.45\linewidth}
\centering
\begin{tikzpicture}
\begin{axis}[otSchema, width=72mm, height=44mm,
  xmin=0, xmax=15, ymin=0.4, ymax=3.6, xtick=\empty, ytick=\empty,
  xlabel={$n$}, xlabel style={at={(1,0)},anchor=west}]
  \draw[otNavy,dashed,line width=0.5pt] (axis cs:0,2) -- (axis cs:15,2);
  \addplot[only marks,mark=*,mark size=1.3pt,draw=otNavy,fill=white,samples at={1,...,14}]
    {2+1.5/(x+0.5)};
  \addplot[only marks,mark=*,mark size=1.3pt,draw=otNavy,fill=white,samples at={1,...,14}]
    {2-1.5/(x+0.5)};
  \addplot[only marks,mark=*,mark size=1.6pt,otGreen,samples at={1,...,14}]
    {2+1.5*sin(deg(1.3*x))/(x+0.5)};
  \node[anchor=east,font=\small,text=otNavy] at (axis cs:0,2) {$\ell$};
  \node[anchor=west,font=\scriptsize,text=otNavy] at (axis cs:14.2,2.3) {$w_n$};
  \node[anchor=west,font=\scriptsize,text=otNavy] at (axis cs:14.2,1.7) {$u_n$};
  \node[anchor=south,font=\scriptsize,text=otGreenDark] at (axis cs:3,3.05) {$v_n$};
\end{axis}
\end{tikzpicture}\\[0.5mm]
\legende{$(v_n)$, prise en étau, est contrainte de converger vers $\ell$.}
\end{minipage}

\begin{theoreme}[Convergence monotone]
Toute suite croissante et majorée converge. Toute suite décroissante et minorée converge.
\end{theoreme}

\begin{propriete}
Une suite croissante non majorée tend vers $+\infty$. Si $(u_n)$ est croissante et converge vers
$\ell$, alors $u_n\leq\ell$ pour tout entier $n$.
\end{propriete}

\begin{attention}
Une suite croissante majorée par $M$ converge vers une limite $\ell\leq M$, mais $\ell$ n'est pas
forcément égale à $M$ : la suite définie par $u_n=1-\dfrac{1}{n+1}$ est majorée par $5$ et converge vers $1$.
\end{attention}

\begin{demonstration}[du programme : limite de $(q^n)$ lorsque $q>1$]
\etape{Inégalité de Bernoulli} Pour tout réel $a>0$ et tout entier naturel $n$, $(1+a)^{n}\geq1+na$.
Pour $n=0$ : $1\geq1$. Si $(1+a)^{k}\geq1+ka$, on multiplie par $1+a>0$ :
$(1+a)^{k+1}\geq(1+ka)(1+a)=1+(k+1)a+ka^{2}\geq1+(k+1)a$.\par\vspace{1mm}
\etape{Conséquence} Si $q>1$, on pose $a=q-1>0$. Alors $q^{n}=(1+a)^{n}\geq1+na$.
Comme $\displaystyle\lim_{n\to+\infty}(1+na)=+\infty$, le théorème de comparaison donne
$\displaystyle\lim_{n\to+\infty}q^{n}=+\infty$.
\end{demonstration}

\essSection{5}{Suites définies par une relation $u_{n+1}=f(u_n)$}

\noindent
\begin{minipage}[c]{0.5\linewidth}
On visualise les premiers termes à l'aide de la courbe de $f$ et de la droite d'équation $y=x$ :
\begin{enumerate}[label=\otLabel{\arabic*.},leftmargin=1.6em,itemsep=0.8mm,topsep=1.5mm]
  \item placer $u_0$ sur l'axe des abscisses ;
  \item rejoindre verticalement la courbe : on lit $u_1=f(u_0)$ en ordonnée ;
  \item rejoindre horizontalement la droite $y=x$ pour reporter $u_1$ en abscisse ;
  \item recommencer pour obtenir $u_2$, $u_3$…
\end{enumerate}
\vspace{1mm}
La « toile d'araignée » obtenue permet de conjecturer le sens de variation et la limite de la suite.
\end{minipage}\hfill
\begin{minipage}[c]{0.45\linewidth}
\centering
\begin{tikzpicture}
\begin{axis}[otSchema, width=66mm, height=62mm,
  xmin=0, xmax=3, ymin=0, ymax=3, xtick={1,2}, ytick={1,2},
  xlabel={$x$}, ylabel={$y$},
  xlabel style={at={(1,0)},anchor=west}, ylabel style={at={(0,1)},anchor=south,rotate=-90}]
  \addplot[otRule!80!otMuted,line width=0.6pt,domain=0:2.9] {x};
  \addplot[otNavy,line width=1.1pt,domain=0:2.9] {0.5*x+1};
  \addplot[otGreen,line width=0.8pt] coordinates
    {(0,0) (0,1) (1,1) (1,1.5) (1.5,1.5) (1.5,1.75) (1.75,1.75) (1.75,1.875) (1.875,1.875)};
  \draw[otGreen!70!otMuted,densely dotted,line width=0.5pt] (axis cs:1,1) -- (axis cs:1,0);
  \draw[otGreen!70!otMuted,densely dotted,line width=0.5pt] (axis cs:1.5,1.5) -- (axis cs:1.5,0);
  \draw[otGreen!70!otMuted,densely dotted,line width=0.5pt] (axis cs:1.75,1.75) -- (axis cs:1.75,0);
  \addplot[only marks,mark=*,mark size=1.3pt,otGreen] coordinates {(0,0) (1,0) (1.5,0) (1.75,0)};
  \node[anchor=north,font=\scriptsize] at (axis cs:0,-0.05) {$u_0$};
  \node[anchor=north,font=\scriptsize] at (axis cs:1,-0.05) {$u_1$};
  \node[anchor=north,font=\scriptsize] at (axis cs:1.5,-0.05) {$u_2$};
  \node[anchor=north,font=\scriptsize] at (axis cs:1.78,-0.05) {$u_3$};
  \fill[otNavy] (axis cs:2,2) circle (1.4pt);
  \node[anchor=south east,font=\scriptsize,text=otNavy] at (axis cs:2.9,2.35) {$\Cf$};
  \node[anchor=north west,font=\scriptsize,text=otMuted] at (axis cs:2.5,2.4) {$y=x$};
\end{axis}
\end{tikzpicture}\\[0.5mm]
\legende{$f(x)=\frac{1}{2}x+1$ et $u_0=0$ : la suite croît vers $2$.}
\end{minipage}

\begin{savoirfaire}[Étudier une suite $u_{n+1}=f(u_n)$]
\begin{enumerate}[label=\otLabel{\arabic*.},leftmargin=1.6em,itemsep=0.6mm,topsep=0.5mm]
  \item Encadrer $u_n$ par récurrence, en utilisant souvent le sens de variation de $f$.
  \item Étudier le sens de variation : signe de $u_{n+1}-u_n$, ou récurrence sur $u_n\leq u_{n+1}$.
  \item Conclure à la convergence grâce au théorème de convergence monotone.
  \item Déterminer la limite $\ell$ : lorsque $f$ est continue, elle vérifie $f(\ell)=\ell$ (voir chapitre 4).
\end{enumerate}
\end{savoirfaire}

\essSection{6}{Algorithme de seuil}

Pour une suite croissante qui tend vers $+\infty$, on cherche le plus petit entier $n$ tel que
$u_n>A$. Exemple avec $u_0=3$ et $u_{n+1}=1{,}5\,u_n-1$ :

\begin{python}
def seuil(A):
    n = 0
    u = 3                  # premier terme u0
    while u <= A:          # tant que le seuil n'est pas dépassé
        n = n + 1
        u = 1.5 * u - 1    # relation de récurrence
    return n
\end{python}

\begin{attention}[Pièges fréquents]
\begin{itemize}[label={\color{otOrange}\textbullet},leftmargin=1.4em,itemsep=0.6mm,topsep=0.5mm]
  \item Une forme indéterminée n'est pas un résultat : il faut transformer l'expression.
  \item Une suite croissante ne converge pas forcément : elle doit aussi être majorée.
  \item « Tendre vers $+\infty$ », ce n'est pas converger : une telle suite est divergente.
  \item Dans une somme géométrique, vérifier que $q\neq1$ avant d'appliquer la formule.
\end{itemize}
\end{attention}

%% =====================================================================
%%  EXERCICES
%% =====================================================================
\exercices
\rubrique{Découverte}

\begin{exercice}[1]{Limites usuelles}
Déterminer la limite de chacune des suites suivantes, définies pour tout entier $n\geq1$.
\begin{tasks}(3)
  \task $u_n=\dfrac{1}{n}$
  \task $v_n=\dfrac{1}{n^{2}}$
  \task $w_n=n^{3}$
  \task $t_n=\sqrt{n}$
  \task $s_n=5+\dfrac{2}{n}$
\end{tasks}
\end{exercice}

\begin{exercice}[1]{Opérations sur les limites}
Déterminer les limites suivantes.
\begin{tasks}(2)
  \task $\displaystyle\lim_{n\to+\infty}(n+7)$
  \task $\displaystyle\lim_{n\to+\infty}\left(3n\times\frac{1}{n}\right)$
  \task $\displaystyle\lim_{n\to+\infty}\frac{n+1}{n}$
  \task $\displaystyle\lim_{n\to+\infty}\left(2-\frac{5}{n^{2}}\right)$
\end{tasks}
\end{exercice}

\rubrique{Entraînement}

\begin{exercice}[2]{Suites géométriques}
Déterminer, si elle existe, la limite de chacune des suites suivantes.
\begin{questions}
  \item $(u_n)$ est définie par $u_0=2$ et, pour tout entier $n$, $u_{n+1}=3u_n$.
  \item $(v_n)$ est définie par $v_0=5$ et, pour tout entier $n$, $v_{n+1}=\frac{1}{2}v_n$.
  \item $(w_n)$ est définie, pour tout entier $n$, par $w_n=4\times\left(\frac{1}{3}\right)^{n}$.
  \item $(t_n)$ est définie, pour tout entier $n$, par $t_n=(-2)^{n}$.
\end{questions}
\end{exercice}

\begin{exercice}[2]{Formes indéterminées}
Lever l'indétermination, puis déterminer la limite.
\begin{tasks}(2)
  \task $\displaystyle\lim_{n\to+\infty}(n^{2}-n)$
  \task $\displaystyle\lim_{n\to+\infty}\frac{3n^{2}+2n+1}{5n^{2}-4}$
  \task $\displaystyle\lim_{n\to+\infty}\frac{2n^{3}-n+1}{n^{2}+3n}$
  \task $\displaystyle\lim_{n\to+\infty}(n-\sqrt{n})$
  \task $\displaystyle\lim_{n\to+\infty}\frac{4n+3}{2n^{3}-n^{2}+1}$
  \task $\displaystyle\lim_{n\to+\infty}\left(\sqrt{n^{2}+n}-n\right)$
  \task $\displaystyle\lim_{n\to+\infty}\frac{n^{4}-3n^{2}+2}{n^{4}+5n^{3}-n}$
\end{tasks}
\end{exercice}

\begin{exercice}[2]{Théorème des gendarmes (1)}
À l'aide d'un encadrement, déterminer les limites suivantes.
\begin{tasks}(2)
  \task $\displaystyle\lim_{n\to+\infty}\frac{\sin n}{n}$
  \task $\displaystyle\lim_{n\to+\infty}\frac{n^{2}+\cos n}{n^{2}+1}$
  \task $\displaystyle\lim_{n\to+\infty}\frac{(-1)^{n}}{n}$
  \task $\displaystyle\lim_{n\to+\infty}\frac{3n+2\sin n}{n+5}$
  \task $\displaystyle\lim_{n\to+\infty}\frac{n^{3}+\cos n}{n^{3}-2}$
  \task $\displaystyle\lim_{n\to+\infty}\frac{(-1)^{n}+n}{n^{2}}$
  \task $\displaystyle\lim_{n\to+\infty}\frac{5n^{2}-1+3\sin n}{5n^{2}+2}$
\end{tasks}
\end{exercice}

\begin{exercice}[2]{Théorème des gendarmes (2)}
Même consigne, avec des encadrements plus délicats.
\begin{tasks}(2)
  \task $\displaystyle\lim_{n\to+\infty}\frac{n+(-1)^{n}\sqrt{n}}{n+1}$
  \task $\displaystyle\lim_{n\to+\infty}\frac{2n^{2}+n\cos n}{2n^{2}-n+1}$
  \task $\displaystyle\lim_{n\to+\infty}\frac{\sin n+\cos n}{n}$
  \task $\displaystyle\lim_{n\to+\infty}\frac{n^{4}+(-1)^{n}n^{2}}{n^{4}+3n}$
  \task $\displaystyle\lim_{n\to+\infty}\frac{4n+2\sin n-1}{4n+3}$
  \task $\displaystyle\lim_{n\to+\infty}\frac{n^{2}-n+\cos n}{n^{2}+2}$
\end{tasks}
\end{exercice}

\begin{exercice}[2]{Récurrence et sommes}
Démontrer par récurrence les égalités suivantes.
\begin{questions}
  \item Pour tout entier $n\geq1$ : $1+2+\dots+n=\dfrac{n(n+1)}{2}$.
  \item Pour tout entier $n\geq1$ : $1^{2}+2^{2}+\dots+n^{2}=\dfrac{n(n+1)(2n+1)}{6}$.
  \item Pour tout entier naturel $n$ : $\displaystyle\sum_{k=0}^{n}2^{k}=2^{n+1}-1$.
  \item Pour tout entier $n\geq1$ : $\dfrac{1}{1\times2}+\dfrac{1}{2\times3}+\dots+\dfrac{1}{n(n+1)}=\dfrac{n}{n+1}$.
\end{questions}
\end{exercice}

\begin{exercice}[2]{Récurrence, inégalités et divisibilité}
Démontrer par récurrence les propriétés suivantes.
\begin{questions}
  \item Pour tout entier naturel $n$ : $2^{n}\geq n+1$.
  \item Pour tout entier $n\geq1$ : $3^{n}\geq2n^{2}+1$.
  \item Pour tout entier naturel $n$ : $4^{n}-1$ est divisible par $3$.
  \item Pour tout entier $n\geq4$ : $2^{n}\geq n^{2}$.
\end{questions}
\end{exercice}

\begin{exercice}[3]{Suites arithmético-géométriques}
Pour chacune des suites suivantes, définies pour tout entier naturel $n$, exprimer le terme général
en fonction de $n$, puis étudier la convergence.
\begin{questions}
  \item $u_0=1$ et $u_{n+1}=2u_n+3$.
  \item $v_0=5$ et $v_{n+1}=\frac{1}{2}v_n+2$.
  \item $w_0=0$ et $w_{n+1}=3w_n-6$.
  \item $t_0=10$ et $t_{n+1}=0{,}8\,t_n+4$.
  \item $s_0=3$ et $s_{n+1}=-\frac{1}{3}s_n+2$.
\end{questions}
\end{exercice}

\rubrique{Problèmes}

\begin{exercice}[3]{Une suite récurrente avec une racine carrée}
On considère la suite $(u_n)$ définie par $u_0=1$ et, pour tout entier naturel $n$,
$u_{n+1}=\sqrt{u_n+2}$.
\begin{questions}
  \item Démontrer par récurrence que, pour tout entier naturel $n$, $0\leq u_n\leq2$.
  \item Montrer que, pour tout entier naturel $n$,
    $u_{n+1}-u_n=\dfrac{(2-u_n)(1+u_n)}{\sqrt{u_n+2}+u_n}$, puis en déduire que $(u_n)$ est croissante.
  \item En déduire que la suite $(u_n)$ converge.
  \item On admet que sa limite $\ell$ vérifie $\ell=\sqrt{\ell+2}$. Déterminer $\ell$.
\end{questions}
\end{exercice}

\begin{exercice}[3]{Devoir de synthèse}
\partieExo{A}{Limites et opérations}
\begin{questions}
  \item Déterminer $\displaystyle\lim_{n\to+\infty}\frac{2n^{2}+3n-1}{n^{2}+5}$.
  \item Déterminer $\displaystyle\lim_{n\to+\infty}\left(n^{3}-2n^{2}\right)$.
  \item Déterminer la limite de la suite définie par $u_n=3\times\left(\frac{2}{3}\right)^{n}$.
\end{questions}
\partieExo{B}{Théorème des gendarmes}
\begin{questions}[resume]
  \item Déterminer $\displaystyle\lim_{n\to+\infty}\frac{2n+\sin n}{n+1}$.
  \item Déterminer $\displaystyle\lim_{n\to+\infty}\frac{n^{2}-1}{n^{2}+\cos n}$.
\end{questions}
\partieExo{C}{Une suite récurrente}
On considère la suite $(v_n)$ définie par $v_0=0$ et, pour tout entier naturel $n$, $v_{n+1}=\dfrac{v_n+3}{4}$.
\begin{questions}[resume]
  \item Montrer que, pour tout entier naturel $n$, $0\leq v_n\leq1$.
  \item Montrer que la suite $(v_n)$ est croissante.
  \item En déduire que $(v_n)$ converge, puis déterminer sa limite.
\end{questions}
\end{exercice}

\begin{exercice}[3]{Une suite homographique}
On considère la suite $(u_n)$ définie par $u_0=\dfrac{1}{3}$ et, pour tout entier naturel $n$,
$u_{n+1}=\dfrac{2u_n}{1+u_n}$.
\partieExo{A}{Étude de la suite $(u_n)$}
\begin{questions}
  \item \begin{questions}
      \item Calculer $u_1$ et $u_2$.
      \item Démontrer par récurrence que, pour tout entier naturel $n$, $0<u_n<1$.
    \end{questions}
  \item \begin{questions}
      \item Montrer que, pour tout entier naturel $n$, $u_{n+1}-u_n=\dfrac{u_n(1-u_n)}{1+u_n}$.
      \item En déduire que la suite $(u_n)$ est croissante.
    \end{questions}
\end{questions}
\partieExo{B}{Une suite auxiliaire}
Pour tout entier naturel $n$, on pose $v_n=\dfrac{u_n}{1-u_n}$, qui est bien défini d'après la partie A.
\begin{questions}[resume]
  \item \begin{questions}
      \item Montrer que, pour tout entier naturel $n$, $v_{n+1}=2v_n$.
      \item En déduire la nature de la suite $(v_n)$, sa raison et son premier terme.
    \end{questions}
  \item \begin{questions}
      \item Exprimer $v_n$ en fonction de $n$.
      \item En déduire que, pour tout entier naturel $n$, $u_n=\dfrac{2^{n}}{2^{n}+2}$.
    \end{questions}
\end{questions}
\partieExo{C}{Limite}
\begin{questions}[resume]
  \item Déterminer la limite de la suite $(v_n)$, puis celle de la suite $(u_n)$.
  \item Vérifier que cette limite $\ell$ est solution de l'équation $x=\dfrac{2x}{1+x}$. Que représente cette équation pour la suite $(u_n)$ ?
\end{questions}
\end{exercice}

\begin{exercice}[3]{Une suite qui converge très vite}
On considère la fonction $f$ définie sur $\intFF{0}{1}$ par $f(x)=x(2-x)$, et la suite $(u_n)$
définie par $u_0=\dfrac{1}{2}$ et, pour tout entier naturel $n$, $u_{n+1}=f(u_n)$.

\noindent
\begin{minipage}[c]{0.56\linewidth}
\partieExo{A}{Conjecture graphique}
\begin{questions}
  \item Calculer les valeurs exactes de $u_1$, $u_2$ et $u_3$.
  \item Sur la figure ci-contre, construire $u_1$, $u_2$ et $u_3$ sur l'axe des abscisses.
    Quelle conjecture peut-on faire sur le sens de variation et la limite de $(u_n)$ ?
\end{questions}
\partieExo{B}{Convergence}
\begin{questions}[resume]
  \item Étudier les variations de $f$ sur $\intFF{0}{1}$.
  \item Démontrer par récurrence que, pour tout entier naturel $n$, $0<u_n<u_{n+1}<1$.
  \item En déduire que la suite $(u_n)$ converge.
\end{questions}
\end{minipage}\hfill
\begin{minipage}[c]{0.4\linewidth}
\centering
\begin{tikzpicture}
\begin{axis}[otGraphe, width=58mm, xmin=0, xmax=1.08, ymin=0, ymax=1.08,
  xtick={0.5,1}, ytick={0.5,1}, xticklabels={,}, yticklabels={$\frac12$,$1$},
  minor tick num=4, grid=both, minor grid style={otRule!45,line width=0.2pt}]
  \addplot[otRule!80!otMuted,line width=0.6pt,domain=0:1.05] {x};
  \addplot[otNavy,line width=1.1pt,domain=0:1,samples=80] {x*(2-x)};
  \addplot[only marks,mark=*,mark size=1.4pt,otGreen] coordinates {(0.5,0)};
  \node[anchor=north,font=\scriptsize] at (axis cs:0.5,-0.05) {$u_0$};
  \node[anchor=south east,font=\scriptsize,text=otNavy] at (axis cs:0.3,0.62) {$\Cf$};
  \node[anchor=north west,font=\scriptsize,text=otMuted] at (axis cs:0.72,0.6) {$y=x$};
\end{axis}
\end{tikzpicture}
\end{minipage}

\partieExo{C}{Expression explicite et vitesse de convergence}
Pour tout entier naturel $n$, on pose $v_n=1-u_n$.
\begin{questions}[start=6]
  \item Montrer que, pour tout entier naturel $n$, $v_{n+1}=v_n^{2}$.
  \item Démontrer par récurrence que, pour tout entier naturel $n$, $v_n=\left(\dfrac{1}{2}\right)^{2^{n}}$.
  \item En utilisant l'inégalité $2^{n}\geq n$, en déduire la limite de la suite $(u_n)$.
  \item On considère la fonction Python ci-dessous. Que renvoie \lstinline[style=otPython]!seuil(9)! ? Interpréter.
\end{questions}
\begin{python}
def seuil(p):
    n = 0
    v = 0.5               # v0 = 1 - u0
    while v >= 10**(-p):
        n = n + 1
        v = v**2
    return n
\end{python}
\end{exercice}

\begin{exercice}[3]{Les abonnés d'une plateforme}
Une plateforme de vidéo à la demande compte $50\,000$ abonnés au 1\textsuperscript{er} janvier 2026.
Une étude montre que, chaque année, 15\,\% des abonnés résilient leur abonnement et que $12\,000$
nouveaux abonnés s'inscrivent. On modélise le nombre d'abonnés, en milliers, au 1\textsuperscript{er} janvier
de l'année $2026+n$ par une suite $(u_n)$, avec $u_0=50$.
\partieExo{A}{Modélisation}
\begin{questions}
  \item Calculer $u_1$ et $u_2$, puis interpréter $u_2$ dans le contexte.
  \item Justifier que, pour tout entier naturel $n$, $u_{n+1}=0{,}85\,u_n+12$.
  \item Pour tout entier naturel $n$, on pose $v_n=u_n-80$. Montrer que $(v_n)$ est géométrique ;
    préciser sa raison et son premier terme.
  \item En déduire que, pour tout entier naturel $n$, $u_n=80-30\times0{,}85^{n}$.
  \item Étudier le sens de variation de $(u_n)$, puis sa limite. Interpréter.
\end{questions}
\partieExo{B}{Objectif 75\,000 abonnés}
\begin{questions}[resume]
  \item Recopier et compléter la fonction suivante, pour qu'elle renvoie l'année à partir de laquelle
    la plateforme compte plus de $75\,000$ abonnés.
\begin{python}
def annee_objectif():
    n = 0
    u = 50
    while u <= ...:
        n = n + 1
        u = ...
    return 2026 + n
\end{python}
  \item Quelle valeur renvoie cette fonction ? On pourra s'aider de la calculatrice.
\end{questions}
\partieExo{C}{Chiffre d'affaires}
Un abonnement annuel coûte 120\,€ ; chaque personne abonnée au 1\textsuperscript{er} janvier paie l'année entière.
\begin{questions}[resume]
  \item Pour tout entier naturel $n$, exprimer $S_n=u_0+u_1+\dots+u_n$ en fonction de $n$.
  \item En déduire le chiffre d'affaires total réalisé de 2026 à 2035 inclus, arrondi au millier d'euros.
\end{questions}
\end{exercice}

\begin{exercice}[3]{Le problème de Bâle}
Pour tout entier $n\geq1$, on pose
\[
  u_n=\sum_{k=1}^{n}\frac{1}{k^{2}}=1+\frac{1}{4}+\frac{1}{9}+\dots+\frac{1}{n^{2}}
  \qquad\text{et}\qquad w_n=u_n+\frac{1}{n}.
\]
\partieExo{A}{Convergence de $(u_n)$}
\begin{questions}
  \item Calculer $u_1$, $u_2$ et $u_3$.
  \item Montrer que la suite $(u_n)$ est croissante.
  \item \begin{questions}
      \item Vérifier que, pour tout entier $k\geq1$, $\dfrac{1}{(k+1)^{2}}\leq\dfrac{1}{k}-\dfrac{1}{k+1}$.
      \item Démontrer par récurrence que, pour tout entier $n\geq1$, $u_n\leq2-\dfrac{1}{n}$.
    \end{questions}
  \item En déduire que la suite $(u_n)$ converge vers un réel $\ell\leq2$.
\end{questions}
\partieExo{B}{Encadrement de $\ell$}
\begin{questions}[resume]
  \item Montrer que, pour tout entier $n\geq1$, $w_{n+1}-w_n=-\dfrac{1}{n(n+1)^{2}}$.
    En déduire le sens de variation de $(w_n)$.
  \item Justifier que, pour tout entier $n\geq1$, $u_n\leq\ell\leq w_n$.
  \item Compléter la fonction Python ci-dessous pour qu'elle renvoie $u_n$, puis, à l'aide de
    \lstinline[style=otPython]!u(1000)!, donner un encadrement de $\ell$ d'amplitude $10^{-3}$.
\begin{python}
def u(n):
    s = 0
    for k in range(1, n + 1):
        s = s + ...
    return s
\end{python}
  \item En 1735, Leonhard Euler a démontré que $\ell=\dfrac{\pi^{2}}{6}$.
    Vérifier que ce résultat est cohérent avec l'encadrement obtenu.
\end{questions}
\end{exercice}
